\section{Finite tabulation and the first crossing number}
\label{sec:tabulation}

\begin{proposition}[A finite list of diagram codes]
For $c\ge1$, all oriented one-component planar diagrams with $c$
crossings, up to subdivided planar redraw, admit at most
\[
                         (c+2)4^c c!
\]
marked codes. Consequently there are finitely many diagrammatic knot
types admitting at most $n$ crossings. More precisely one may bound their
number by $1+\sum_{c=1}^{n}(c+2)4^cc!$.
\end{proposition}
\begin{proof}
Start just before any crossing visit, traverse the knot, and number visits
$1,\ldots,2c$. The Dowker even--odd rule proved above pairs the odd visits
bijectively with the even visits. Record that permutation of $c$ symbols.
On each pair record which visit goes over and the sign of the crossing.
This gives four choices per pair. There are $c+2$ possible exterior
faces on the connected sphere map, by Euler's formula. Mark the exterior
face by its index when face orbits are sorted by least dart.

The traversal tells which leaving dart joins which next arriving dart.
At a crossing of visits $i,j$, their incoming and outgoing darts are
opposite in pairs. Of the two possible oriented cyclic orders of these
four named darts, precisely one has the recorded crossing sign for the
recorded overpass. Thus the code recovers the oriented rotation table,
the edge pairing, and the overpass and orientation data. Successive
cyclic predecessors (or successors with the other fixed convention)
after edge reversal trace the faces. Discard codes with an invalid
permutation, face index, or Euler count. The argument only requires
that every diagram is represented. Distinct codes may describe the same
knot.

Here is why subdivision coordinates add no multiplicity to a fixed code.
Take disjoint small polygonal disks about the crossings and narrow strips
along the intervening arcs. Their cell patterns are prescribed by the
rotation table. Subdivide strip lengths on both drawings to match, and
triangulate their rectangular pieces in matching order, keeping the
diagram as center lines. The remaining bounded regions are disks with
simple polygon boundaries: thickening a spanning tree gives a disk and
inserting each remaining edge splits the corresponding complementary
disk in the sphere. For the marked exterior face use a large bounding
polygon; the extra region between its boundary and the outside boundary
of the thickening is an annulus. A narrow strip from one matched marked
point of the inner boundary to the outside box cuts this annulus to a
disk. Such a polygonal path can be drawn along the exterior collar to an
extreme vertex and then out along a supporting ray. The pieces can all
be taken rational: at each disjointness requirement take a positive
rational clearance smaller than the finitely many relevant separations.
For disjoint rational segments the squared separation is computed by
testing the endpoint-to-segment projections, clamped to the segment,
in both directions. For a squared separation $q>0$, half of
$\min(1,q)$ is a valid strict clearance. Direction and angular margins
at a vertex use the finite list of nonzero rational determinants.

Two polygonal disks with corresponding boundary vertices have compatible
triangulated subdivisions. To see this, triangulate each polygon by
diagonals. A short constructive diagonal argument is as follows. Remove
collinear subdivision points temporarily. Choose a strictly convex
vertex $v$ (one maximizing a generic linear functional suffices), with
neighbors $u,w$. If the closed triangle $uvw$ contains no further
vertex, $uw$ cuts off an ear, unless the polygon is already a triangle.
Otherwise join $v$ to a further vertex of this triangle farthest from
the line $uw$. The open segment so obtained has greater height above
$uw$ than any such further vertex. An unrelated polygon edge could only
enter the triangle through its base or end in the triangle, since it
cannot cross $vu,vw$. Hence it cannot cross that open segment.
The segment starts inside at the convex corner, and so is an interior
diagonal. Repeat on the smaller polygons and reinsert all boundary
subdivision points. Transport each triangulation to the same strictly
convex rational polygon with the matched boundary order (points
$(i,i^2)$ in increasing order, closed by the chord between the extreme
points, suffice). There the two triangulations overlay in convex cells;
triangulate those cells and pull the refinements back by the respective
rational affine maps on triangles. This gives the common triangulations
required for the redraw certificate. Together the matched disks, strips,
and exterior pieces give the claimed redraw.

Enumerating permutations, bit choices, and marked faces is bounded; a
coarser $c^c$-entry scan also suffices to enumerate the permutations for
$c\ge1$. The bound $c!$ follows by successively recording the rank of the
partner among the remaining even positions. The arithmetic statement of
finiteness is the finite coding and collision principle: in any supplied
list of knot diagrams with the indicated crossing cap and length greater
than the bound, two have identical positive-crossing codes and a redraw
witness, or two have zero crossings. Pigeonhole follows by counting
filled slots in a bounded Boolean array. Every oriented simple polygon
redraws to a triangle of matching orientation by the disk and exterior
construction. The two orientations of the unknot have a short move
witness. Orient the bow-tie polygon with successive vertices
$(-1,-1),(1,1),(-1,1),(1,-1)$ and choose either overpass at its double
point. Its upper and lower triangular lobes have opposite orientations.
Deleting one empty lobe or the other gives the two orientations of a
simple polygon. This justifies the one zero-crossing slot for knot types.
\end{proof}

\begin{theorem}[The smallest crossing number of a nontrivial knot is three]
Every knot diagram having one or two crossings admits a sequence of
type I deletions and planar redraws to a crossing-free diagram. The
trefoil realizes three and is nontrivial.
\end{theorem}
\begin{proof}
At two crossings the cyclic word of visits uses letters $A,B$ twice
each. The interlaced order $ABAB$ would pair two odd visits, contrary
to the even--odd rule. Up to cyclic shift and relabeling the word is
therefore $AABB$. The subarc between consecutive visits to $A$, closed
at its crossing, is a simple polygonal lobe. The analogous $B$-lobe is
disjoint from it. At the base of a lobe its two rays belong to different
passages and are adjacent in the cyclic order. The remainder of the knot
away from that base is a connected arc, so it lies entirely on one side
of the lobe. Choose a lobe whose bounded interior does not contain the
other lobe, by choosing the inner one if they are nested. Its interior
then misses the whole remainder of the diagram. It is an empty monogon.
A thin disk neighborhood of the lobe together with its empty interior,
extended slightly past the crossing, therefore has exactly the type I
curl pattern; the overpass choice permits one of the two type I
deletions. The polygon triangulation construction above supplies the local
redraw when the lobe has many subdivision vertices.

At a single crossing there are two such lobes, now meeting only at
their common base. Away from the base each lies entirely on one side
of the other. At least one bounded interior is empty of the other lobe:
if one contains the other away from their common point, its own boundary
away from that point is exterior to the contained lobe. This is also the
disk separation test on the two simple polygons. Delete the empty lobe
as before. The resulting crossing-free diagram represents the unknot.
The closed positive braid $\sigma_1^3$ has three crossings and Jones
polynomial $t+t^3-t^4\ne1$, as calculated in
Section~\ref{sec:bracket}. This proves both bounds.
\end{proof}
